\begin{theorem}[G-23: Exact eventual multiplicity sign in every strict direction through $m$ rows]\label{G-23}
Let $n>m\ge17$ and $4\le r\le m$. Put $a=n-m$, and fix
$\nu\vdash ns$, $s\ge1$, with
\[
 \nu_r>0,\qquad \nu_1>as.
\]
Append no additional positive parts beyond row $r$.
Define
\[
 \begin{split}
 \eta&=\nu-\nu_r(1,\ldots,1),\qquad e=|\eta|,\qquad
 r_\eta=\#\{i<j:\eta_i>\eta_j\},\\
 u&=(r-2)n^2+r_\eta-r^2+3,\\
 v&=rm^2-2m+r+r_\eta-r^2+2,\\
 \Delta&=v-u=rm^2-(r-2)n^2-2m+r-1.
 \end{split}
\]
The integer $\Delta$ is never zero. Both multiplicities are
positive at every sufficiently large integer stretch, and
\begin{equation}\label{G23:ratio}
 \frac{m_{j\nu}(\ell^a\operatorname{per}_m)}
      {m_{j\nu}(\det_n)}
 =\kappa_\nu j^\Delta(1+O(j^{-1})),\qquad \kappa_\nu>0,
\end{equation}
through all integers $j\to\infty$, in original coordinate degree $sj$.
Thus the permanent multiplicity is eventually strictly larger
exactly when
\begin{equation}\label{G23:boundary}
 n\le N_{m,r}:=
 \left\lfloor\sqrt{\left\lfloor
       \frac{rm^2-2m+r-1}{r-2}\right\rfloor}\right\rfloor;
\end{equation}
for $n>N_{m,r}$ the determinant multiplicity is eventually strictly larger.

Here are explicit bounds. Let
\[
 \begin{aligned}
 B_n&=\binom{n+r-1}{r-1},&
 b&=\binom{B_n+s-1}{s}\dim S_\eta\mathbb C^r,\\
 L&=ns+(r-1)e,&p&=ns+e,\\
 \beta&=\left\lceil\frac{3r^2b}{8}\right\rceil
                          L^{2(r^2-1)},&
 T&=\binom{b+\beta}{b}-1,\qquad E=rT(\beta!).
 \end{aligned}
\]
The Schur dimension is the Weyl product
$\prod_{i<j}(\eta_i-\eta_j+j-i)/(j-i)$.
Both multiplicities are positive for every $j\ge\beta^2$.
If $\Delta>0$, the strict positive gap holds for every
\[
 j\ge E\bigl(4u(pE)^{u-1}+1\bigr).
\]
If $\Delta<0$, the strict reverse gap holds for every
\[
 j\ge E\bigl(4v(pE)^{v-1}+1\bigr).
\]
\end{theorem}

\paragraph{Proof dependencies and meaning.}
G-21 gives the permanent image dimension, R-28 gives the
determinant dimension and the projectively free common stable
tensor, and G-22 supplies stability for every stated direction.
The invariant-generator input is Derksen (2001), Theorem~1.1
and Proposition~1.2, printed p.~956. The Hilbert-function and
rational image-degree inputs are G-10 and its exact source
locators. The changes from G-18 are the growing row count,
the different group dimension, the new stable-point construction,
and the uniform arithmetic exclusion of equality. We give the
argument with all new parameters, rather than assuming that
the four-row constants remain valid.

\begin{proof}
Use the two enhanced cones of G-22, and let $A_D,A_P$ be
their invariant coordinate rings. They are graded domains.
Their degree-$j$ dimensions are the two specified multiplicities.
By R-28 and G-21,
\[
 \dim\mathcal D_{n,r}=(r-2)n^2+2,\qquad
 \dim\mathcal P_{m,r}=rm^2-2m+2.
\]
The padded permanent cone has dimension
$rm^2-2m+r+1$. Indeed its projective parametrization by a
padding line and a permanent form is generically injective:
on the nonempty open where the latter is irreducible of degree
$m>1$ and not divisible by the padding line, the unique
nonlinear factor and the linear factor recover both projective
inputs. The irreducible open is nonempty by the common cyclic
family in R-28. Thus padding adds $r-1$ affine-cone dimensions.
Subtracting $r^2-1$ from the enhanced dimensions in G-22 proves
\[
 \dim A_D=u,\qquad\dim A_P=v.
\]
Both are at least two: $r\le m<n$ gives
$u\ge2n^2-r^2+3>2$, and
$v\ge(r-1)m^2-2m+r+2>2$.

\paragraph{1. Invariant degrees and absence of a period.}
The ambient representation has dimension $b$. Its action
coefficients on $\mathrm{SL}_r$ have degree at most $L$:
the covariant factor contributes $ns$, and the entries of
the inverse transpose have degree $r-1$, contributing at
most $(r-1)e$ on the Schur factor. Its kernel is finite
because it is contained in the finite stabilizer of the
common stable tensor. In the embedding
$\mathrm{SL}_r\subset\mathbb A^{r^2}$, the defining equation
has degree $r$ and the group dimension is $r^2-1$.
Derksen's Proposition~1.2 therefore gives
\[
 \sigma(W)\le rL^{r^2-1},
\]
and Theorem~1.1 gives the displayed integer generator bound
$\beta$. Exactness of invariants makes $A_D,A_P$ graded
quotients of this same ambient invariant ring.

Choose homogeneous generators $g_i$ of degrees
$1\le e_i\le\beta$. R-28 supplies a common stable point
$w$ with
\[
 gw=cw,\quad g\in\mathrm{SL}_r,\ c\in\mathbb C^*
                  \quad\Longrightarrow\quad c=1.
\]
For precision, its projective stabilizer is the center
$\mu_r$, and this center acts with exponent
$ns-e=r\nu_r$, hence trivially.
The degrees of the generators nonzero at $w$ have greatest
common divisor one. If that divisor were $h>1$, a primitive
$h$th root $c$ would make every generator have equal values
at $w$ and $cw$, although their two closed orbits are distinct.
This contradicts closed-orbit separation by invariants
(Mukai, Theorem~5.3 and Corollary~5.4, pp.~161--162).
There is at least one such generator by the same separation
applied to the closed orbits of $w$ and zero.

The numerical-semigroup proof in G-18, part~2, now applies
with this actual $\beta$: shortest paths modulo the least
degree have at most $\beta-1$ edges, so each residue has a
representative of weight less than $\beta^2$. Hence every
integer $j\ge\beta^2$ is a sum of degrees nonzero at $w$.
The corresponding product gives an ambient homogeneous
invariant $h_j$ nonzero on both cones. This proves both
occurrence assertions at every such integer $j$.

\paragraph{2. Standard Veroneses and effective comparison.}
There are at most $T$ chosen positive-degree generators,
because the whole ambient coefficient-polynomial space
through degree $\beta$ has dimension $T$.
The generator-grouping argument of G-11 applies with
$E=rT\beta!$: remove pure generator blocks of degree
$\beta!$, leaving a residual monomial of degree less than
$T\beta!$, then complete it to a block of degree $E$ and
partition the remaining pure blocks. Thus $A_D^{(E)}$ and
$A_P^{(E)}$ are standard graded domains. Since $T\ge\beta$
and $\beta!\ge\beta$, we have $E\ge\beta^2$.

Both enhanced cones have dense homogeneous parametrizations
of degree $p$ in independent variables. They are
\[
 (Z,B)\longmapsto\det\left(\sum_{i=1}^r Z_i x_i\right)^{\otimes s}
                  \otimes S_\eta(B)v_\eta
\]
and
\[
 (Z,H,B)\longmapsto
 \left(H^a\operatorname{per}_m\left(\sum_{i=1}^rZ_i x_i\right)\right)^{\otimes s}
                  \otimes S_\eta(B)v_\eta.
\]
The parameter $H$ is an independent linear form and $B$ an
independent endomorphism of $V^*$. The first factors have
parameter degree $ns$, and the second has degree $e$.
If $e=0$ it is constant and may be omitted. These are parameter
degrees, distinct from the action degree $L$.
G-10's base-locus-aware degree argument bounds the two
projective invariant-Veronese image degrees by
\[
 (pE)^{u-1}\quad\hbox{and}\quad(pE)^{v-1},
\]
respectively.

Write $j=Ek+t$, $0\le t<E$. For either domain $A$, multiplication
by the common nonzero invariants $h_{E+t}$ and $h_{2E-t}$ gives
\begin{equation}\label{G23:sandwich}
 \dim A_{E(k-1)}\le\dim A_j\le\dim A_{E(k+2)}\qquad(k\ge1).
\end{equation}
If $v>u$, linear Noether normalization of $A_P^{(E)}$ and
the Chardin bound on $A_D^{(E)}$ yield
\[
 \dim(A_P)_{E(k-1)}\ge\binom{k+u-1}{u},\qquad
 \dim(A_D)_{E(k+2)}\le(pE)^{u-1}\binom{k+u+1}{u-1}.
\]
Their binomial ratio is
\[
 \frac{k(k+1)(k+2)}{u(k+u)(k+u+1)}
                 \ge\frac{k}{4u}\qquad(k\ge u+1).
\]
The asserted positive threshold follows from
\eqref{G23:sandwich}. If $u>v$, exchange the domains and
replace $u$ by the actual dimension $v$ in this argument.
This gives the asserted reverse threshold without an upper
bound substituted for the permanent dimension.

For a standard graded domain of dimension $q\ge2$, the
leading Hilbert term is
$\deg(\operatorname{Proj} A^{(E)})k^{q-1}/(q-1)!$,
with error $O(k^{q-2})$. Applying \eqref{G23:sandwich}
shows, through all integers, that
\[
 \dim A_j=c_A j^{q-1}+O(j^{q-2}),\qquad c_A>0.
\]
Taking the ratio for $A_P,A_D$ proves \eqref{G23:ratio}.

\paragraph{3. The integer dimensions can never be equal.}
If $r$ is even, $\Delta$ is odd. Suppose $r$ is odd and
$\Delta=0$. Reduction modulo two forces $m,n$ to have
the same parity. Whether both are even or both are odd,
reduction modulo four then gives
\[
 0=\Delta\equiv r-1\pmod4,
\]
so $r\equiv1\pmod4$. Reduction modulo $r-2$ gives
$2m^2-2m+1\equiv0$, and hence
\[
 (2m-1)^2\equiv-1\pmod{r-2}.
\]
But $r-2\equiv3\pmod4$, so one of its prime divisors
$p$ is congruent to three modulo four. A square root of
$-1$ modulo $p$ is impossible: its nonzero residue would
have $(p-1)$st power both $1$, by Fermat's theorem, and
$(-1)^{(p-1)/2}=-1$. This contradiction proves
$\Delta\ne0$. The precise integer boundary
\eqref{G23:boundary} follows immediately.
\end{proof}

\paragraph{Rectangular consequence and limits.}
For a rectangular direction, the strict condition is
$n<rm/(r-1)$. Throughout that range the exponent is positive:
\[
 \Delta>
 \frac{r m^2}{(r-1)^2}-2m+r-1
 =\frac r{(r-1)^2}\left(m-\frac{(r-1)^2}{r}\right)^2
                                      +\frac{r-1}{r}>0.
\]
Thus every such rectangle gives a gap at every sufficiently
large admissible coordinate degree. For a fixed $m,n,r$,
one may take $s=r/\gcd(n,r)$ and
$\nu=(ns/r)^r$, so all admissible degrees are $sj$.
The theorem concerns fixed strict directions with $r\le m$.
It does not decide the face $\nu_1=as$, the range $r>m$,
every earlier finite degree, or a uniform polynomial-time
positive rule when the parameters grow. Its bounds and
non-occurrence-obstruction conclusion retain these limits.

